\SourceLeaf{279}%
Berechnung mit zehnstelligen Logarithmen, vollständig ausreicht. Wer diese Tafel zur Berechnung von Messungen innerhalb dieser Zone verwenden will, kann nach Belieben einige der letzten Dezimalstellen weglassen, wenn ihm eine geringere Genauigkeit genügt (und das ist der häufigste Fall). In welcher Form man die \textit{Resultate} einer Messung auch darstellt, sollte man dies konsequenterweise immer mit einer Genauigkeit tun, die der Genauigkeit der Messungen selbst entspricht. Aus den Zahlen der Resultate sollte man die \mbox{beobachteten} Größen immer wieder so genau zurückgewinnen können, wie sie gemessen wurden. Wählt man dafür also ausschließlich Längen und Breiten, so würde man selbst trigonometrischen Messungen von nur mäßiger Genauigkeit keineswegs gerecht, wenn man die Resultate nur mit der Genauigkeit angäbe, mit der sich Längen und Breiten auf astronomischem Weg bestimmen lassen. Dadurch erhielte man lediglich einen falschen Maßstab für die Güte der Arbeit und brächte sich häufig gerade um die aussagekräftigsten Prüfungen dieser Güte.

\Article{11}
\indent Die hier betrachtete konforme Abbildung der Ellipsoidfläche auf die Kugelfläche lässt sich auf mehr als eine Weise zur Berechnung trigonometrischer Messungen nutzen. In dieser Abhandlung geht es nur um die unmittelbare Anwendung. Andere, daraus abgeleitete Anwendungen für denselben Zweck bleiben einer zweiten Abhandlung vorbehalten.

Die unmittelbare Anwendung ist im Wesentlichen bereits in der oben genannten Schrift kurz umrissen. Ein aus kürzesten Linien, den sogenannten geodätischen Linien, gebildetes Dreieckssystem auf der Ellipsoidfläche wird auf der Kugelfläche durch ein Dreieckssystem dargestellt, dessen Winkel den entsprechenden Winkeln auf dem Sphäroid genau gleich sind. Sofern die Seiten keine Meridianbögen sind, weichen sie im strengen Sinn von Großkreisbögen ab, aber so wenig, dass man sie in den meisten Fällen als völlig mit ihnen zusammenfallend betrachten darf. Wo höchste Genauigkeit verlangt wird, lässt sich die Abweichung zumindest leicht mit der nötigen Genauigkeit berechnen. Dabei ist stets vorausgesetzt, dass
